\newpage
\chapter{Memoria Descriptiva}
\label{chap:memoria}

El presente proyecto consiste en el diseño de un sistema de cableado estructurado
para la Escuela Profesional de Enfermería de la Universidad Nacional de San Antonio
Abad del Cusco (UNSAAC), con la finalidad de implementar una infraestructura de
telecomunicaciones moderna, segura y escalable que permita soportar los servicios
de conectividad institucional.

La infraestructura propuesta contempla una red convergente capaz de integrar
servicios de transmisión de datos, acceso inalámbrico, telefonía IP y futuras
aplicaciones tecnológicas como sistemas de videovigilancia IP. El diseño se basa
en estándares internacionales de cableado estructurado utilizando cableado UTP
Categoría 6 para el subsistema horizontal y fibra óptica para el backbone principal.

Asimismo, la red estará organizada mediante un Main Distribution Frame (MDF)
ubicado en el cuarto nivel del edificio, desde donde se distribuirán los enlaces
hacia las áreas académicas y administrativas.

El proyecto considera además criterios de ordenamiento, etiquetado,
administración y mantenimiento de la infraestructura, garantizando confiabilidad,
facilidad de expansión y adecuado desempeño de los servicios de telecomunicaciones.

% -----------------------------------------------------------
\section{Descripción general del proyecto}
\label{sec:descripcion_general}
% -----------------------------------------------------------

El proyecto consiste en el diseño e implementación de un sistema de cableado
estructurado categoría 6 para el edificio de la Escuela Profesional de Enfermería,
identificado institucionalmente como dentro de la red de
conectividad de la UNSAAC.

El sistema contempla la instalación de \textbf{29 puntos de red} distribuidos en
tres niveles del edificio (primero, segundo y cuarto), todos ellos terminados en
el rack principal \textbf{G09 (MDF)}, ubicado en el cuarto nivel. Desde este punto,
el edificio se integra a la red institucional de la UNSAAC mediante un enlace de
fibra óptica monomodo OS2 tendido a través de la canalización subterránea del campus,
administrada por la Dirección de Tecnologías de la Información y Comunicaciones (DTIC).

El diseño abarca el subsistema horizontal en cable UTP categoría 6, el subsistema
de backbone en fibra óptica, las canalizaciones (canaletas PVC, conduit EMT y
bandejas portacables), la sala de telecomunicaciones con su rack de 42U, el sistema
de puesta a tierra y la red inalámbrica (WLAN) con siete Access Points Wi-Fi 6.

% -----------------------------------------------------------
\section{Objetivos de la obra}
\label{sec:objetivos}
% -----------------------------------------------------------

\subsection{Objetivo general}

Diseñar e implementar un sistema de cableado estructurado categoría 6 en la
Escuela Profesional de Enfermería de la UNSAAC, que provea una infraestructura
de telecomunicaciones confiable, ordenada y escalable, integrada a la red
institucional del campus universitario.

\subsection{Objetivos específicos}

\begin{itemize}
    \item Dotar a todos los ambientes académicos y administrativos del edificio
    de puntos de red certificados bajo la norma \mbox{ANSI/TIA-568-C.2}.

    \item Establecer un backbone de fibra óptica que interconecte el MDF del
    edificio con la red de campus de la DTIC-UNSAAC mediante tecnología Gigabit
    y 10 Gigabit Ethernet.

    \item Implementar una sala de telecomunicaciones que centralice
    la gestión del cableado estructurado, conforme a \mbox{ANSI/TIA-569-D} y
    \mbox{ANSI/TIA-606-B}.

    \item Diseñar un sistema de puesta a tierra para telecomunicaciones que
    garantice la equipotencialización de los equipos activos y la protección
    contra interferencias electromagnéticas, según \mbox{ANSI/TIA-607-B}.

    \item Proveer cobertura inalámbrica WiFi 6 en el 97\% del área útil del
    edificio, con niveles de señal de al menos $-70$\,dBm en aulas y oficinas.

    \item Documentar el sistema mediante planos, esquemas y etiquetado completo
    que faciliten la administración y el mantenimiento futuro de la infraestructura.
\end{itemize}

% -----------------------------------------------------------
\section{Ubicación y área de intervención}
\label{sec:ubicacion}
% -----------------------------------------------------------

\subsection{Ubicación geográfica}

\begin{table}[H]
    \centering
    \caption{Datos de ubicación del área de intervención.}
    \label{tab:ubicacion}
    \begin{tabular}{|p{4.5cm}|p{9cm}|}
        \hline
        \textbf{Parámetro} & \textbf{Descripción} \\
        \hline
        Departamento       & Cusco \\
        \hline
        Provincia          & Cusco \\
        \hline
        Distrito           & Cusco \\
        \hline
        Institución        & Universidad Nacional de San Antonio Abad del Cusco
                             (UNSAAC) \\
        \hline
        Facultad           & Enfermería \\
        \hline
        Escuela Profesional & Enfermería \\
        \hline
        Identificación     & CIUDAD UNIVERSITARIA DE PERAYOC UNSAAC \\
        \hline
    \end{tabular}
\end{table}%
\nobreak\vspace{-0.3cm}
\subsection{Descripción del inmueble}

El edificio de la Escuela Profesional de Enfermería es una construcción de
concreto armado y muros de albañilería confinada de cuatro niveles. Cuenta con
una planta típica de aproximadamente 680\,m² por nivel, con una distribución
funcional que incluye ambientes administrativos en el primer nivel, aulas en el
segundo y ambientes académicos y de reunión en el cuarto nivel.

\begin{table}[H]
    \centering
    \caption{Distribución funcional del edificio G09 por nivel.}
    \label{tab:distribucion_niveles}
    \begin{tabular}{|c|p{5cm}|p{6.5cm}|}
        \hline
        \textbf{Nivel} & \textbf{Función principal} & \textbf{Ambientes principales} \\
        \hline
        Primero & Área administrativa y biblioteca &
        Decanato, Secretarías, Jefatura, Sala de Lectura,
        Almacén de Libros, Editorial, A.S.P., Hall \\
        \hline
        Segundo & Área académica &
        Aulas 201--204, Galería (Aula 210),
        Ambiente 209, Hall, SS.HH. \\
        \hline
        Cuarto  & Área académica / MDF &
        Ambientes de trabajo, Ambientes 404--405,
        Salón de Grados, Sala de Telecomunicaciones \\
        \hline
    \end{tabular}
\end{table}%
\nobreak\vspace{-0.3cm}
% -----------------------------------------------------------
\section{Alcances del proyecto}
\label{sec:alcances}
% -----------------------------------------------------------

El proyecto comprende el diseño, suministro e instalación de los siguientes
componentes y subsistemas:

\textbf{Cableado horizontal:} Tendido de cable UTP categoría 6 de cuatro pares
desde el rack G09 hasta cada uno de los 29 puntos de red del edificio, con una
longitud máxima de 90\,m por canal conforme a \mbox{ANSI/TIA-568-C.2}. Cada
punto de red cuenta con un outlet doble de dos puertos RJ-45 Cat.6.

\textbf{Backbone de fibra óptica:} Enlace de fibra óptica monomodo entre
el nodo de distribución y el MDF del edificio, y enlaces
de fibra óptica multimodo OM3 para el backbone vertical interno entre el MDF y
los equipos activos de los niveles inferiores.

\textbf{Sala de telecomunicaciones y rack:} Instalación del Gabinete G09 de 42U
en el cuarto nivel, equipado con panel de distribución de fibra óptica (ODF),
tres patch panels Cat.6 de 24 puertos, switch de distribución de 24 puertos GbE
con uplinks SFP+, UPS de 1500\,VA y sistema de ventilación con termostato.

\textbf{Canalizaciones:} Suministro e instalación de canaletas PVC tipo V0, bandejas portacables metálicas y cajas
de registro, conforme a \mbox{ANSI/TIA-569-D}.

\textbf{Puesta a tierra:} Implementación del sistema de equipotencialización según \mbox{ANSI/TIA-607-B},
con resistencia total verificada.

\textbf{Red inalámbrica (WLAN):} Instalación de siete Access Points Wi-Fi 6 con soporte dual-band, MU-MIMO y OFDMA, alimentados por PoE+
desde el switch del rack G09, cubriendo el 97\% del área útil del edificio con
un nivel de señal mínimo de $-70$\,dBm.

\textbf{Documentación:} Elaboración del expediente técnico completo, que incluye
planos en AutoCAD, memoria descriptiva, especificaciones técnicas, metrados y
presupuesto.

Lo que no forma parte del alcance del presente proyecto: obras civiles
adicionales al edificio existente, instalaciones eléctricas de baja tensión
(más allá de la puesta a tierra de telecomunicaciones), equipos de usuario final
(computadoras, teléfonos IP) ni la configuración lógica de los equipos de red
más allá de la puesta en marcha básica.

% -----------------------------------------------------------
\section{Justificación técnica y beneficios}
\label{sec:justificacion}
% -----------------------------------------------------------

La Escuela Profesional de Enfermería carece actualmente de una infraestructura
de telecomunicaciones certificada y debidamente documentada. El cableado existente,
en muchos casos improvisado, no cumple con los estándares mínimos de desempeño
ni de administración, lo que genera problemas de conectividad, dificultad para
el diagnóstico de fallas y una capacidad de respuesta muy limitada ante el
crecimiento de la demanda tecnológica en la institución.

La implementación del sistema de cableado estructurado categoría 6 resuelve
esta situación, aportando los siguientes beneficios técnicos e institucionales:

\begin{table}[H]
    \centering
    \caption{Beneficios técnicos e institucionales del proyecto.}
    \label{tab:beneficios}
    \begin{tabular}{|p{4.5cm}|p{9cm}|}
        \hline
        \textbf{Beneficio} & \textbf{Descripción} \\
        \hline
        Ancho de banda garantizado &
        El cableado Cat.6 soporta aplicaciones de hasta 1\,GbE por enlace,
        con margen para 10GbE en distancias cortas, suficiente para las
        necesidades actuales y futuras del edificio. \\
        \hline
        Cobertura inalámbrica total &
        Los siete APs Wi-Fi 6 garantizan conectividad sin interrupciones
        para hasta 610 usuarios simultáneos en todo el edificio. \\
        \hline
        Administración simplificada &
        El etiquetado completo y la documentación del sistema permiten
        identificar y resolver fallas en menor tiempo, reduciendo el
        impacto sobre las actividades académicas. \\
        \hline
        Escalabilidad &
        La infraestructura tiene capacidad de crecimiento sin necesidad
        de rediseño, tanto en puntos de red cableados como en APs
        inalámbricos adicionales. \\
        \hline
        Integración institucional &
        El sistema se integra a la red de campus de la UNSAAC mediante
        fibra óptica, facilitando el acceso a servicios centralizados
        de la DTIC (Internet, servidores, VoIP). \\
        \hline
        Seguridad y protección &
        La puesta a tierra normalizada y los protectores de sobretensión
        resguardan los equipos activos frente a fallas eléctricas y
        descargas, prolongando su vida útil. \\
        \hline
    \end{tabular}
\end{table}

Desde el punto de vista académico, el proyecto mejora directamente las condiciones
tecnológicas para la enseñanza, el acceso a plataformas virtuales, la realización
de videoconferencias y el uso de recursos digitales por parte de docentes y
estudiantes de la Escuela Profesional de Enfermería.%
\nobreak\vspace{-0.3cm}
% -----------------------------------------------------------
\section{Normas y estándares aplicados}
\label{sec:normas}
% -----------------------------------------------------------

El diseño e implementación del sistema de cableado estructurado se rige por las
normas internacionales de mayor reconocimiento en la industria de
telecomunicaciones. La Tabla \ref{tab:normas} resume los estándares aplicados
y su ámbito de aplicación en el proyecto.

\begin{table}[H]
    \centering
    \caption{Normas y estándares internacionales aplicados al proyecto.}
    \label{tab:normas}
    \begin{tabular}{|p{4.5cm}|p{9cm}|}
        \hline
        \textbf{Norma / Estándar} & \textbf{Ámbito de aplicación} \\
        \hline
        ANSI/TIA-568-C.0 &
        Requisitos generales de cableado de telecomunicaciones en edificios
        comerciales. Topología, distancias y arquitectura del sistema. \\
        \hline
        ANSI/TIA-568-C.1 &
        Cableado de telecomunicaciones en edificios comerciales.
        Requisitos mínimos de puntos de red y subsistemas. \\
        \hline
        ANSI/TIA-568-C.2 &
        Especificaciones para cableado de par trenzado no blindado (UTP)
        categoría 6. Parámetros de desempeño y pruebas de certificación. \\
        \hline
        ANSI/TIA-568-C.3 &
        Especificaciones para cableado de fibra óptica. Parámetros de
        atenuación, conectores y empalmes. \\
        \hline
        ANSI/TIA-569-D &
        Espacios y sendas para telecomunicaciones en edificios comerciales.
        Diseño de salas de telecomunicaciones, canalizaciones y rutas de
        cableado. \\
        \hline
        ANSI/TIA-606-B &
        Administración de infraestructura de telecomunicaciones.
        Etiquetado, codificación y documentación del sistema. \\
        \hline
        ANSI/TIA-607-B &
        Puesta a tierra y equipotencialización para telecomunicaciones.
        Diseño del sistema TGB, TMGB y conductores de bonding. \\
        \hline
        ISO/IEC 11801:2017 &
        Estándar internacional de cableado genérico para edificios y
        campus. Especificaciones de clases y categorías de cableado,
        fibra óptica y topología de red. \\
        \hline
        IEEE 802.11ax (Wi-Fi 6) &
        Estándar para redes inalámbricas de área local. Tecnologías
        OFDMA, MU-MIMO, dual-band y seguridad WPA3. \\
        \hline
        IEEE 802.3at (PoE+) &
        Alimentación eléctrica sobre cableado Ethernet para dispositivos
        como Access Points y teléfonos IP (hasta 30\,W por puerto). \\
        \hline
        NEC / NFPA 70 &
        Código Eléctrico Nacional. Referencia para la instalación
        eléctrica asociada al sistema de telecomunicaciones. \\
        \hline
        IEEE 81 &
        Guía para la medición de resistencia de puesta a tierra mediante
        telurómetro. \\
        \hline
    \end{tabular}
\end{table}

El cumplimiento de estas normas garantiza que el sistema de cableado estructurado
instalado en la Escuela Profesional de Enfermería sea interoperable con cualquier
tecnología de red activa compatible con los estándares IEEE 802.3 y 802.11,
y que su desempeño pueda ser verificado y certificado de forma objetiva mediante
equipos de medición homologados.